% profiles/michigan-state-university.tex — Michigan State University
% ============================================================================
% source: https://grad.msu.edu/academics/theses-and-dissertations/formatting-your-document
%         (The Graduate School, "Formatting Your Document" — the page that
%          serves the two documents below)
% source: https://edge.sitecorecloud.io/michigansta4a14-msu70a4-prod718c-8cd0/media/project/msu/grad/manuals/etd-formatting-guide-august-2026.pdf
%         ("MSU Formatting Guide for Theses and Dissertations", footed
%          "August 2026" on every page — linked from that page as the
%          "Printable Formatting Guide")
% source: https://edge.sitecorecloud.io/michigansta4a14-msu70a4-prod718c-8cd0/media/project/msu/grad/manuals/sample-page-packet-jan-2026.pdf
%         ("Printable Sample Pages", January 2026; the annotated thesis and
%          dissertation title pages are pages 2 and 3 of 18)
% accessed: 2026-09-09
% checked:  texlib-thesis-profile-round
%
% VERIFIED and set here: geometry (a range — see below), spacing (a choice —
% likewise), the title page, and the absence of a separate approval page.
% NOT set:
%   * \thesissetchapteropening — no deeper chapter-opening margin is prescribed.
%     (The title page's own 2in top margin IS prescribed, and is applied inside
%     \thesistitlepage rather than through that hook, which governs chapters.)
%   * every checkable accessibility declaration — see the accessibility note.
%
% CURRENCY, AND A STALE EDITION THAT SEARCH ENGINES STILL SERVE. The governing
% guide is dated August 2026 and lives on MSU's Sitecore CDN. Two older
% editions are still reachable at the OLD grad.msu.edu/sites/default/files path
% and are what a search for "MSU ETD formatting guide" returns first:
% "ETD Formatting Guide August 2025.pdf" and "ETD Formatting Guide FINAL Jan
% 2023.pdf". Neither is linked from the current Formatting Your Document page.
% Follow the CDN link that page serves; do not take a figure from the
% grad.msu.edu/sites/... copies.
%
% THE APPROVAL IS A FORM, NOT A PAGE. "The signed Approval Form for electronic
% submission serves as evidence that the document has been examined and approved
% by the major professor (or thesis/dissertation director) and guidance
% committee. The electronic approval form is received by The Graduate School
% after the major professor ... has signed off on it." The guide goes further and
% forbids carrying it with the manuscript: "Do not attach the Approval Form,
% Embargo Form, IRB letter or IACUC letter as supplemental files in your ProQuest
% account." Its list of preliminary pages — Title, Abstract, Copyright,
% Dedication, Acknowledgements, Preface, Table of Contents, and the lists — has
% no approval or signature leaf in it. Hence \thesisnoapprovalpage.
%
% SIGNATURES: yes, but never inside the document. The major professor signs the
% electronic Approval Form, which reaches the Graduate School separately.
%
% THE TYPEFACE RULE IS A BAN, NOT A LIST, so this class satisfies it — unlike at
% Arizona State, whose closed list of eight faces it fails. MSU: "Script, small
% caps, and ornamental fonts will not be accepted. Fonts strongly suggested
% include Times New Roman, Calibri, Arial, Garamond, and Cambria." Latin Modern
% is none of the three banned kinds, and the named five are "strongly suggested"
% rather than required. Type size: "A 12-point font is the norm for regular
% text", which the class's 12pt report already is.

\thesisinstitution{Michigan State University}

% "All regular margins must be between 1 and 1 1/2 inches and consistent on all
% sides. The only text permitted to be in the margins are the page numbers."
%
% A RANGE, and the profile takes the bottom of it. 1in is in spec, is what the
% class defaults to, and keeps the most text on the page; anything up to 1.5in
% is equally in spec provided it is used on all four sides. Note the consistency
% clause — a binding allowance on the left alone would be out of spec here.
%
% The page-number relaxation is the reason the class's `includefoot' geometry is
% safe: MSU expects the folio "1/2 inch to an inch from the bottom of the page
% (in the footer)", and permits it in the margin.
\thesissetgeometry{left=1in, right=1in, top=1in, bottom=1in}

% "The manuscript may be 1.5-spaced or double-spaced and must be consistent
% throughout the document." A CHOICE; the profile takes double, which is the
% class default. \thesissetspacing{onehalf} is equally in spec. The guide lists
% the usual single-spaced exceptions: table-of-contents entries within chapters,
% multi-line list entries, "titles, headings, footnotes, endnotes, lengthy
% quotations".
\thesissetspacing{double}

% --- Accessibility -----------------------------------------------------------
% The August 2026 Formatting Guide was read end to end. The string "accessib"
% appears in it exactly ONCE, and not as a document-conformance rule: it is
% about transcribing linked media. There is no WCAG level, no tagging rule, no
% alt-text rule, no document-language rule and no PDF standard anywhere in it,
% and no Graduate School accessibility page for ETDs was located.
%
% MSU does run a university-wide digital accessibility programme, and its
% library publishes guidance on accessible LaTeX. Neither is the Graduate
% School's filing requirements, so neither is a source for this file and neither
% is encoded here.
\thesisaccessibilityrequirement{Not published}
	{The Graduate School's Formatting Guide states no accessibility requirement
	 for the document itself — no WCAG level, no tagging, no alt text, no
	 document language, no PDF standard. It was read in full on the access date.
	 This records that the Graduate School asks for nothing, not that nobody
	 looked.}
\thesisaccessibilityrequirement{Transcribe linked media}
	{``QR codes/YouTube videos/etc.\ will need to be transcribed within the
	 document for accessibility purposes.'' Links to outside resources ``must be
	 maintained in perpetuity'', and for media the student created ``a
	 summary/transcription must be provided'' immediately following it.}

% --- Title page ----------------------------------------------------------------
% Reproduced from the annotated dissertation specimen (page 3 of the Sample
% Pages packet), whose green inch ruler is what fixes the vertical positions,
% together with the guide's own bullets:
%
%   - "Do not use boldface type on the Title Page."
%   - "The title is to be in all capital letters and single-spaced", and the
%     specimen's own text says "CENTERED, 2 INCHES FROM THE TOP OF THE PAGE".
%   - "Capitalize ``By'' and Your Name."
%   - "``A THESIS'' or ``A DISSERTATION'' is to be in all capital letters."
%   - "The text block beneath A THESIS or A DISSERTATION should be centered,
%     left to right, and appear as follows: 1. Submitted to 2. Michigan State
%     University 3. in partial fulfillment of the requirements 4. for the degree
%     of 5. <Graduate Degree Granting Unit/Program> - <degree that is being
%     conferred> 6. <the year of degree conferral>", with the specimen's callout
%     "These 4 lines should be single-spaced" against lines 1--4.
%   - "The year in line 6 must match the year listed in the ProQuest account".
%
% THE SEPARATOR IN LINE 5 IS AMBIGUOUS AND IS LEFT AS THE PROSE WRITES IT. The
% guide's text uses a hyphen-minus, "<unit/program> - <degree>"; the rendered
% specimen shows what looks like an em dash with no space before the degree
% ("Mechanical Engineering ---Doctor of Philosophy"). RESEARCHING.md says to
% prefer the prose where a template and the prose disagree, so a spaced hyphen
% is what is emitted. A reviewer who reads the specimen as authoritative should
% change one character here.
%
% The unit/program string is \thesis@program and must be one of the Graduate
% Degree Granting Units on the Registrar's list: "Only the units/programs listed
% here are approved by Michigan State University."
%
% Vertical placement: the 2in top is a stated figure and is set exactly. The
% three gaps below it are NOT stated anywhere, so they are elastic in the
% proportions the specimen's ruler shows (title to "By" ~1.3in, name to
% "A DISSERTATION" ~2.9in, year to the bottom margin ~0.9in) rather than frozen,
% which keeps a long title from pushing the block off the page.
\renewcommand{\thesistitlepage}{%
	\loadgeometry{nopagenumber}%
	\begin{titlepage}%
		\thispagestyle{empty}\centering\singlespacing
		% geometry's top margin is 1in, so 1in more puts the title at 2in.
		\vspace*{1in}%
		\thesis@tagtitle{Title}{\MakeUppercase{\@title}}\par
		\vspace{\stretch{1.3}}
		By\par
		\vspace{\baselineskip}
		\@author\par
		\vspace{\stretch{2.9}}
		A \MakeUppercase{\thesis@DocNoun}\par
		\vspace{\baselineskip}
		Submitted to\\
		Michigan State University\\
		in partial fulfillment of the requirements\\
		for the degree of\par
		\vspace{\baselineskip}
		\thesis@program\ - \thesis@degree\par
		\vspace{\baselineskip}
		\thesis@date\par
		\vspace{\stretch{0.9}}
	\end{titlepage}%
	\loadgeometry{normal}%
}

% --- Committee approval page ---------------------------------------------------
% None: the approval is an electronic Approval Form that the Graduate School
% receives separately, and the guide's list of preliminary pages contains no
% approval or signature leaf. See the header.
\thesisnoapprovalpage
